\subsection{Finite free norm identities with genuine nonetale multiplicities}
\label{subsec:finite-free-norms}

These statements apply to arbitrary commutative-ring algebras that are finite
free over their base. Nilpotent layers contribute genuine local lengths;
inseparable residue extensions contribute their full degrees. The coefficient
seminorm may be arbitrary. A finite projective determinant and transport of
regularizable decreasing limits remain separate foundations for the complete
finite flat statement in FHX \texttt{lem:finite-flat-norm}.

\begin{theorem}[Weighted local norm over every complete field]
\label{spectrum:nonetale-local-weighted-norm}
\lean{EnvelopeProperty.wgaussValuation_normForm_local_of_algClosureSplit,
EnvelopeProperty.wgaussValuation_normForm_of_local}
\leanok
\uses{algebra:finite-artinian-local-factors,spectrum:algclosed-local-norm-form,
spectrum:whole-norm-form-base-change,spectrum:local-algebra-norm-length,
spectrum:complete-norm-extension}
Let $H$ be a complete ultrametric normed field and $E$ any finite-dimensional
commutative local $H$-algebra. If $w$ on $\kappa(E)$ extends the norm of $H$,
then, for finite families $b_j\in E$ and weights $r_j\geq0$,
\[
 \left\|N_{E/H}\Bigl(\sum_jX_jb_j\Bigr)\right\|_r
 =\left(\max_jw(\overline b_j)r_j\right)^{\ell_E(E)[\kappa(E):H]}.
\]
No separability, reducedness or infinitude is required.
\end{theorem}
\begin{proof}
Choose an extending valuation on an algebraic closure $\Omega$ of $H$.
The actual Artinian localization product of $\Omega\otimes_HE$ consists
of finite local $\Omega$-algebras $L_i$. Whole-polynomial determinant base
change and the genuine product isomorphism identify the original norm form
with the product of these local norm forms. The algebraically closed local
calculation writes each as a power of the linear polynomial obtained from
its actual residue scalar, with exponent $\ell_{L_i}(L_i)$.

Restrict a factor's residue scalar along $E\to\Omega\otimes_HE$ to obtain
an actual homomorphism $E\to\Omega$. Its kernel is maximal, since $E$ is
finite over $H$ and its image is a finite-dimensional domain over $H$.
Locality identifies that kernel with the unique maximal ideal of $E$.
Thus it descends to a genuine $H$-embedding $\kappa(E)\to\Omega$.
Complete-field norm uniqueness identifies all its scalar values with $w$.
Every local linear factor consequently has weighted norm
$\max_jw(\overline b_j)r_j$. The sum of the factor lengths equals
$\dim_H E$: each $\Omega$-local factor has residue degree one, and dimension
is preserved by tensor base change and the actual product. The genuine
local dimension identity gives $\dim_HE=\ell_E(E)[\kappa(E):H]$.
Weighted Gauss multiplicativity completes the argument. The modular helper
accepts the real local-factor isomorphism; the unconditional theorem supplies
it from the constructed Artinian decomposition. No interpolation is used.
\end{proof}

\begin{theorem}[Weighted norm of an arbitrary finite algebra]
\label{spectrum:nonetale-finite-weighted-norm}
\lean{EnvelopeProperty.wgaussValuation_normForm_of_localAlgEquiv,
EnvelopeProperty.wgaussValuation_normForm_of_finite}
\leanok
\uses{spectrum:nonetale-local-weighted-norm,algebra:finite-artinian-local-factors,
spectrum:whole-norm-form-base-change}
Let $E$ be any finite-dimensional commutative algebra over a complete
ultrametric field $H$. For every maximal ideal $m$, prescribe a valuation
$w_m$ on $\kappa(m)$ extending the norm of $H$. Then
\[
 \left\|N_{E/H}\Bigl(\sum_jX_jb_j\Bigr)\right\|_r
 =\prod_{m\in\operatorname{Max}(E)}
 \left(\max_jw_m(b_j)r_j\right)^{\ell_{E_m}(E_m)[\kappa(m):H]}.
\]
\end{theorem}
\begin{proof}
The actual algebra isomorphism $E\simeq\prod_mE_m$ factors its entire norm
polynomial into the local norm polynomials. Every localization is genuinely
finite over $H$, and its actual local residue field is $\kappa(m)$.
Apply the preceding local calculation and weighted multiplicativity.
The helper proves the same statement for any supplied genuine product of
finite local algebras; the final theorem supplies the canonical localization
product. The zero algebra gives the empty product, as its determinant norm
is one in rank zero.
\end{proof}

\begin{definition}[Intrinsic completed-fibre lengths and residual valuations]
\label{spectrum:nonetale-fibre-length}
\lean{EnvelopeProperty.Berkovich.Spectrum.fibreLocalLengthENat,
EnvelopeProperty.Berkovich.Spectrum.fibreLocalLength,
EnvelopeProperty.Berkovich.Spectrum.fibreNormMult,
EnvelopeProperty.Berkovich.Spectrum.fibreLocalLength_eq,
EnvelopeProperty.Berkovich.Spectrum.fibreResidualHom,
EnvelopeProperty.Berkovich.Spectrum.fibreResidualHom_algebraMap,
EnvelopeProperty.Berkovich.Spectrum.fibreResidualValuation,
EnvelopeProperty.Berkovich.Spectrum.fibreResidualValuation_algebraMap,
EnvelopeProperty.Berkovich.Spectrum.fibreResidualValuation_one_tmul}
\leanok
\uses{spectrum:fibre-tensor-evaluation,spectrum:etale-fibre-multiplicities}
For $A\to B$, $y\in\mathcal S(B)$ over $x$, let
$E_x=H(x)\otimes_AB$ and let $p_y$ be the actual kernel of $E_x\to H(y)$.
The intrinsic extended local length is $\ell_{(E_x)_{p_y}}((E_x)_{p_y})$.
Its natural version uses $\operatorname{ENat.toNat}$, and the natural raw
multiplicity is that number times the existing actual residue degree.
The genuine residue-field lift $\kappa(p_y)\to H(y)$ defines a valuation
extending the norm of $H(x)$ and taking value $|b|_y$ on $1\otimes b$.
For an arbitrary nonfinite map no geometric norm multiplicity or finite
length is asserted by these definitions.
\end{definition}
\begin{proof}
The tensor map is a real ring homomorphism to the field $H(y)$, so its
kernel is prime and defines an actual support localization. Outside this
kernel the map is nonzero, hence has unit image. The universal residue-field
lift therefore exists and agrees with the tensor map on every original
element. Pull back the norm valuation of $H(y)$. The genuine isometric map
$H(x)\to H(y)$ proves its coefficient compatibility, and the whole tensor
formula on $1\otimes b$ proves its value at every $b$. Naming a different
but equal restricted base point preserves the actual local length.
\end{proof}

\begin{theorem}[Finite local multiplicities, positivity and rank]
\label{spectrum:nonetale-fibre-multiplicities}
\lean{EnvelopeProperty.Berkovich.Spectrum.fibreNormMult_eq,
EnvelopeProperty.Berkovich.Spectrum.fibreNormMult_eq_localFinrank,
EnvelopeProperty.Berkovich.Spectrum.fibreLocalLengthENat_ne_top,
EnvelopeProperty.Berkovich.Spectrum.one_le_fibreNormMult,
EnvelopeProperty.Berkovich.Spectrum.sum_fibreNormMult}
\leanok
\uses{spectrum:nonetale-fibre-length,spectrum:finite-completed-fibres,
spectrum:local-algebra-norm-length,algebra:finite-artinian-local-factors}
For a finite algebra $B/A$, the intrinsic local lengths at actual fibre
points are finite. The genuine norm multiplicity is
\[
 m_y=\ell_{(E_x)_{p_y}}((E_x)_{p_y})[\kappa(p_y):H(x)]
     =\dim_{H(x)}(E_x)_{p_y}\geq1.
\]
If $B/A$ is finite free, then $\sum_{y\text{ over }x}m_y=\operatorname{rank}_AB$.
\end{theorem}
\begin{proof}
Integrality makes $p_y$ maximal. Its quotient is a field and is genuinely
isomorphic to its residue field, giving the stated degree. For a finite
algebra, each localization factor is finite-dimensional over $H(x)$.
It is Artinian and Noetherian as a module over itself, so its extended length
is finite. The local dimension identity identifies the multiplicity with
its dimension. This localization at a prime is nontrivial, giving positivity.
The actual maximal-kernel fibre equivalence reindexes all local factors by
all points over $x$. Sum their dimensions using the real Artinian product,
and use finite free scalar-extension rank to obtain the original rank.
\end{proof}

\begin{theorem}[Scalar finite free norm at every seminorm point]
\label{spectrum:finite-free-scalar-norm}
\lean{EnvelopeProperty.Berkovich.Spectrum.val_algebraNorm_of_finiteFree}
\leanok
\uses{spectrum:nonetale-fibre-multiplicities,spectrum:nonetale-local-valued-norm,
spectrum:algebra-norm-base-change}
For $B/A$ finite free, $x\in\mathcal S(A)$ and $b\in B$,
\[
 |N_{B/A}(b)|_x=\prod_{y\text{ over }x}|b|_y^{m_y}.
\]
\end{theorem}
\begin{proof}
Actual determinant base change evaluates the original norm in $H(x)$
as the norm of $1\otimes b$ in $E_x$. The genuine Artinian localization
product factors that norm. For every maximal factor, use the actual
residue-field map to the completed field of its corresponding fibre point.
Its pulled-back valuation extends the norm of $H(x)$ and takes the value
$|b|_y$. The complete-field scalar local norm calculation retains the true
local length and full residue degree, including inseparable extensions.
Reindex through the genuine maximal-kernel fibre bijection and identify each
exponent with $m_y$. No etaleness or reducedness is required.
\end{proof}

\begin{theorem}[Whole weighted finite free norm at every seminorm point]
\label{spectrum:finite-free-weighted-norm}
\lean{EnvelopeProperty.Berkovich.Spectrum.wgaussValuation_map_normForm_of_finiteFree,
EnvelopeProperty.Berkovich.Spectrum.wgaussValuation_normForm_of_finiteFree}
\leanok
\uses{spectrum:nonetale-finite-weighted-norm,spectrum:nonetale-fibre-multiplicities,
spectrum:whole-norm-form-base-change}
Write the whole determinant norm form as
\[
N_{B/A}\Bigl(\sum_jX_jb_j\Bigr)=\sum_\alpha c_\alpha X^\alpha.
\]
For every $x\in\mathcal S(A)$ and weights $r_j\geq0$,
\[
 \max_\alpha |c_\alpha|_x r^\alpha
 =\prod_{y\text{ over }x}\left(\max_j|b_j|_y r_j\right)^{m_y}.
\]
The original coefficient formula holds without separability, etaleness,
reducedness or infinitude assumptions.
\end{theorem}
\begin{proof}
Whole-polynomial determinant base change identifies the image of the original
norm form in $H(x)[X]$ with the actual norm form of $E_x$.
Use the unconditional finite-algebra weighted theorem with the real residual
valuations of the maximal factors. Each value on $1\otimes b_j$ is
$|b_j|_y$, and each local-length times residue-degree exponent is $m_y$.
Reindex all factors by the actual fibre bijection. Coefficient valuation
preservation along $A\to H(x)$ returns the original weighted coefficient
formula. This is an identity for the actual determinant polynomial; it does
not use evaluation interpolation on $H(x)$.
\end{proof}

\begin{theorem}[Nonetale norm products under actual base change]
\label{spectrum:finite-free-norm-products-base-change}
\lean{EnvelopeProperty.Berkovich.Spectrum.weightedFibreNormProduct_baseChange,
EnvelopeProperty.Berkovich.Spectrum.scalarFibreNormProduct_baseChange,
EnvelopeProperty.Berkovich.Spectrum.weightedFibreNormProduct_scalarExtension,
EnvelopeProperty.Berkovich.Spectrum.scalarFibreNormProduct_scalarExtension,
EnvelopeProperty.Berkovich.Spectrum.weightedFibreNormProduct_relativeScalarExtension,
EnvelopeProperty.Berkovich.Spectrum.scalarFibreNormProduct_relativeScalarExtension}
\leanok
\uses{spectrum:finite-free-weighted-norm,spectrum:finite-free-scalar-norm,
spectrum:whole-norm-form-base-change,spectrum:scalar-extension-algebra-cancellation}
For an arbitrary coefficient algebra $C/A$ and $x'\in\mathcal S(C)$, the
weighted and scalar norm products for $C\otimes_AB/C$ at $x'$ equal those
for $B/A$ at the restricted point. Each fibre uses its own genuine local
length times residue degree. The same comparison holds for the actual
scalar-extended map $S\otimes_R A\to S\otimes_R B$, and for compatible
prescribed coefficient valuations under their true restriction map.
\end{theorem}
\begin{proof}
Replace both weighted products by their actual determinant norm polynomial
values. Whole-polynomial coefficient base change and preservation of the
coefficient seminorm identify those values. For scalar products, replace
each with its actual scalar norm and apply determinant scalar base change.
For $S\otimes_R A\to S\otimes_R B$, use the constructed tensor cancellation
isomorphism and its real whole norm-form compatibility; its finite free
instances follow from the actual transported scalar-extended module.
The prescribed-base versions use the true compatible restriction map and
the same absolute formulas. Fibres need not be in bijection, and no such
bijection is assumed. These are finite free polynomial norm comparisons,
not a claim about all function classes or transport of decreasing limits.
\end{proof}
